% Generated draft: requires mathematical and source review.
\documentclass[11pt]{article}
\usepackage[margin=25mm]{geometry}
\usepackage{fontspec}
\setmainfont{Latin Modern Roman}
\usepackage{amsmath,amssymb,mathtools,unicode-math}
\setmathfont{Latin Modern Math}
\usepackage{enumitem,needspace,xurl,hyperref}
\setlength{\parindent}{0pt}
\setlength{\parskip}{5pt}
\setlength{\emergencystretch}{3em}
\newcommand{\N}{\mathbb{N}}
\newcommand{\Z}{\mathbb{Z}}
\newcommand{\Q}{\mathbb{Q}}
\newcommand{\R}{\mathbb{R}}
\newcommand{\C}{\mathbb{C}}
\newcommand{\F}{\mathbb{F}}
\newcommand{\A}{\mathbb{A}}
\newcommand{\PP}{\mathbb{P}}
\newcommand{\E}{\mathbb{E}}
\newcommand{\eps}{\varepsilon}
\newcommand{\dd}{\,\mathrm{d}}
\begin{document}
\section*{UnsolvedMath ID: 701739}
\textbf{Source problem code:} catalogue entry
\subsection{Definitions and mathematical statement}
Let \(\mathfrak X\) be a class of finite groups. A subgroup is an \(\mathfrak X\)-subgroup if, as a group, it belongs to \(\mathfrak X\). For a subset \(S\) of a group, write \(\langle S\rangle\) for the subgroup generated by \(S\).

The question attributed to D. O. Revin asks whether there exists such a class \(\mathfrak X\) satisfying all three conditions below.
\begin{enumerate}
\item For every \(H\in\mathfrak X\), every subgroup \(K\le H\) belongs to \(\mathfrak X\), and every homomorphic image of \(H\) belongs to \(\mathfrak X\).
\item For every group \(G\), without any finiteness assumption on \(G\), and every pair of normal subgroups \(N,M\trianglelefteq G\),
\[
N,M\in\mathfrak X\quad\Longrightarrow\quad NM\in\mathfrak X,
\qquad NM=\{nm:n\in N,\ m\in M\}.
\]
Here \(N\) and \(M\) are finite because \(\mathfrak X\) consists of finite groups.
\item For every integer \(n\ge1\), there exist a finite group \(G\) and an element \(g\in G\) whose conjugacy class
\[
D=g^G=\{x^{-1}gx:x\in G\}
\]
satisfies
\[
\langle d_1,\ldots,d_n\rangle\in\mathfrak X
\quad\text{for every }(d_1,\ldots,d_n)\in D^n,
\qquad
\langle D\rangle\notin\mathfrak X.
\]
The witnesses \(G\) and \(D\) may depend on \(n\).
\end{enumerate}

\textit{Qualification concerning condition (3).} The displayed formulation interprets the source's phrase ``any \(n\) elements'' as allowing repetitions, equivalently requiring the assertion for every nonempty subset of \(D\) of cardinality at most \(n\). The source does not explicitly specify whether the elements must be distinct. Under an exactly \(n\) distinct elements interpretation, the corresponding requirement would instead be
\[
\forall S\subseteq D,\quad |S|=n\ \Longrightarrow\ \langle S\rangle\in\mathfrak X,
\]
which is vacuous when \(|D|<n\). The intended convention requires confirmation.

\textit{Supplied background qualification.} The record reports that no class satisfying (1)--(3) can also be closed under extensions, attributing this result to D. O. Revin (2025). Here extension closure means that, for every finite group \(E\) and every \(K\trianglelefteq E\),
\[
K\in\mathfrak X\ \text{ and }\ E/K\in\mathfrak X
\quad\Longrightarrow\quad E\in\mathfrak X.
\]
This is a supplied literature assertion, not a result established or independently verified in this draft.
\subsection{Short English statement}
Does there exist a class of finite groups closed under subgroups, homomorphic images, and products of two normal members inside any group, but with the following property? For every positive integer \(n\), some conjugacy class in a finite group has all subgroups generated by \(n\) of its elements in the class, while the subgroup generated by the whole conjugacy class lies outside it. The meaning of ``\(n\) elements'' needs confirmation as described above. The source labels the problem open and reports that extension closure rules out such a class.
\subsection{Sources}
\begin{itemize}
\item Supplied local UnsolvedMath registry record: permanent ID \(701739\), exact source code \texttt{catalogue entry}, title \emph{Kourovka Notebook Problem 20.81}. The supplied status is \texttt{open}; this label has not been independently verified. The record was created and updated on 2026-10-16.
\item \emph{The Kourovka Notebook (Unsolved Problems in Group Theory)}, problem 20.81, proposed by D. O. Revin in 2022. The supplied background identifies the 20th issue (2022) and the 21st issue edition (2026), arXiv:1401.0300v48. The registry describes the collection as edited by E. I. Khukhro and V. D. Mazurov. These source and edition details are supplied attribution and have not been independently checked.
\item D. O. Revin, \emph{Algebra i Analiz}, volume 37, no. 1 (2025), pp. 141--176 (Russian). The record cites this reference for the asserted nonexistence of an extension-closed class satisfying (1)--(3). No article title or external link is supplied, and the cited assertion has not been independently verified.
\end{itemize}
\end{document}
